\section{From operating inputs to the linked statements}
\label{sec:company-input-authoring}

Suppose Northstar's selling price changes from 10 to 11. A stored annual
projection can display a new revenue figure, but that alone does not specify
the customer claim, collection, tax or closing entry that should change with
it. The operating-input successor makes those dependencies explicit. Its
\code{CompanyYearInputsDraft/1} body imports the historical opening and
authors the year's operating quantities and contract terms. Forecast journals
are calculated from those inputs.

The purchase policy first combines demand of 30 units, target inventory of
8 units and opening inventory of 4 units. It purchases 34 units at the
explicit homogeneous cost of 5, producing inventory and a payable of 170.
Delivery of 30 units at price 10 produces a customer claim and revenue of
300, with inventory cost of 150. The stated control-transfer assumption is
required before that revenue entry can be proposed. An unknown control fact
leaves the operation unresolved; a known failure of that condition does not
silently use the ordinary sale method.

Claims are generated by those recognition events. The authored collection
fraction of 90 percent settles 270 of the new claim, while a separately
identified collection settles 10 of opening receivables. Changing price to
11 raises the new claim by 30 and collections by 27. At the demonstration's
immediate 20 percent tax rate, tax rises by 6. Profit therefore rises by 24,
cash by 21 and closing receivables by 3. The difference between profit and
cash follows from the unpaid part of the increased sales, not a balancing
adjustment. Changing demand instead changes purchases and cost of sales
through their quantity rules.

Asset memory is also explicit. The opening plant cost of 122 comprises
cohorts costing 100, 10 and 12, with useful lives of 10, 5 and 12 years.
Their opening accumulated depreciation is 30, 1 and zero. All are in service
at the start of the example year, so their annual charges total 13 under
the stated straight-line, zero-residual method. The new expenditure of 20
enters service on 1 January 2026 and contributes no 2025 depreciation. If the
unused 12-cost cohort also enters service in 2026, the 2025 charge falls by
1. Profit rises by 0.8 while cash tax rises by 0.2. The restricted method
refuses a midyear service date until an accrual fraction is specified; it
does not replace an unavailable convention with a full-year charge.

\paragraph{A shared input body.}
Text and typed records carry the same opening reference, source digest,
context, fiscal period, inventory policy, claims, asset cohorts, debt mix,
facility, tax method and event instants. The closed schema requires every
field. Exact amounts use reduced rational strings, and supported unknowns
retain named obligations. The evaluator verifies the opening's portable
record identity as well as the supplied source bytes. Its permitted
information mode is archival exogenous stress: the example does not pretend
that its historical opening was available before publication.

Each conditional event records its method, input dependencies, proposed
journal and economic-state changes. Claim balances, inventory units and
asset cohorts reconcile with their control accounts after every event. The
tax entry is derived from accrued income accounts and closing transfers
their net balance to retained earnings. An explicit cash check precedes
each payment. The draw and repayment retain their gross financing amounts
even though their net principal effect is zero.

The 18 retained rate/share/borrowing-cost paths reproduce every journal,
cash prefix, closing account and all 29 report lines of the independently
retained comparison, including the additional interest and tax effects.
Removing all forecast-projection records from the imported package leaves
the new calculation unchanged. These checks establish the input-to-result
meaning for this example; changing price, demand or asset service gives
different results through the stated operation rules.

\paragraph{An interrupted year has a partial report.}
When payroll is moved before collections, the available cash is zero and
the payment of 50 fails. The earlier draw and supplier payment remain in
the conditional state; no payroll expense or cash payment is applied. A
partial report identifies that last successful event and its dated balances.
Opening cash and equity remain known, but requested annual cash, profit
and the other uncomputed year-end lines remain unresolved. An unknown
collection instant supplies no invented completed prefix. Complete reports
bind every statement to the same input body and fiscal period, and reject
an attempt to combine a changed-price input with an old result.

The exact schema, both authoring forms, reference semantics and checked
comparisons are retained under
\path{data/accounting-master/company-inputs-20261023/}. This successor
uses one occurrence of each of fourteen event kinds in a full calendar
year, synthetic treatment admission and the stated simple methods. Repeated
invoices, partial-period accruals, revisions, wider contracts and sector
methods need their own definitions. Historical reconstruction, general
source-to-method admission and durable posting remain separate obligations.
The result is conditional and uncommitted.
